%% raw_roc.tex -- raw transcription of the reading "ROC Curves Deep Dive"
%% (source/reading/roc_tutorial.html, page title "ROC Curves: An Overview").
%% Trimmed + compressed: passages already covered in Part I (d3_evaluation, d4_rcode) are omitted.
\section{Reading --- ROC Curves Deep Dive}
\ixh{ROC curve}\ixh{AUC}\ixch{titanic}%
{\color{figgray}\itshape [HTML title ``ROC Curves: An Overview'', PSTAT 131/231. Omitted (in Part I): data loading,
factor, \texttt{coin\_flip}, 70--30 stratified split, Recipe B, workflows, training accuracy (0.807 / 0.616), test re-fit,
confusion-matrix heatmaps (A at 0.5 and 0.00131; B at 0.5 and 0.4076525) with TPR/FPR arithmetic and points,
TPR/FPR definitions, Model B AUC (0.496) and accuracy, \texttt{roc\_curve()} needs probabilities,
AUC = discrimination, two uses of ROC, AUC $=1\Rightarrow$ leakage.]}

\subsection{Recipe A}
\begin{rcode}
titanic_recipe_a <- recipe(survived ~ pclass + sex + age +   (*@\ixch{recipe()}@*)
                           sib_sp + parch + fare, titanic_train) %>%
  step_impute_linear(age, impute_with = imp_vars(sib_sp)) %>%  (*@\ixch{step\_impute\_linear()}\ixch{imp\_vars()}\ixch{impute\_with}\ixh{linear imputation}@*)
  # choice of predictors to impute with is up to you
  step_dummy(all_nominal_predictors()) %>%    (*@\ixch{step\_dummy()}\ixch{all\_nominal\_predictors()}\ixh{dummy variable}@*)
  step_interact(~ starts_with("sex"):age + age:fare)  (*@\ixch{step\_interact()}\ixch{starts\_with()}\ixh{interaction term}@*)
\end{rcode}

\subsection{Model A: ROC curve}
\figref{25}
\begin{routput}
# A tibble: 254 x 3
   .threshold specificity sensitivity      (*@\ixch{.threshold}@*)
 1 -Inf            0                1
 2    0.00131      0                1
 3    0.00474      0.0121           1
 4    0.00563      0.0182           1
 5    0.00635      0.0242           1
 6    0.0116       0.0303           1
 7    0.0117       0.0364           1
 8    0.0122       0.0424           1
 9    0.0147       0.0485           1
10    0.0186       0.0545           1
# ... with 244 more rows
\end{routput}
By default \ixc{roc\_curve()} tries \rc{threshold}\ixh{threshold}{} values in a grid from \rc{-Inf} ($\approx0$) to \rc{Inf} ($\approx1$).
First predicted class probabilities\ixh{posterior probability}:
\begin{rcode}
log_a_results %>% select(.pred_Yes, survived) %>% head()  (*@\ixch{select()}\ixch{.pred\_Yes}\ixch{head()}@*)
\end{rcode}
\begin{routput}
  .pred_Yes survived
1   0.0672  No
2   0.268   No
3   0.792   Yes
4   0.00474 No
5   0.247   Yes
6   0.693   Yes
\end{routput}
Default \ix{threshold} for binary classification\ixh{classification}{} is generally 0.50: No if $p<.50$, else Yes:
\begin{rcode}
augment(log_test_a, new_data = titanic_test) %>%   (*@\ixch{augment()}@*)
  mutate(.pred_class = factor(if_else(.pred_Yes < 0.5, "No", "Yes"), levels = c("Yes", "No"))) %>%  (*@\ixch{mutate()}\ixch{if\_else()}\ixch{factor()}\ixch{.pred\_class}@*)
  conf_mat(truth = survived, estimate = .pred_class) %>%  (*@\ixch{conf\_mat()}@*)
  autoplot(type = "heatmap")
# that point on the ROC curve:
log_a_results %>% roc_curve(survived, .pred_Yes) %>% autoplot() +
  geom_point(aes(x = 0.1212121, y = 0.776699), colour = spec_colors[3], size = 2.5)  (*@\ixch{geom\_point()}\ixch{aes()}@*)
\end{rcode}
Same for 0.25, 0.75, \ldots{} any value in $[0,1]$. Goal: maximize \ix{sensitivity} (TPR\ixh{TPR}) while minimizing
1 $-$ specificity\ixh{specificity}{} (FPR\ixh{FPR}).
\begin{rcode}
log_a_results %>%
  ggplot(aes(x = .pred_Yes, fill = survived)) +   (*@\ixch{ggplot()}\ixch{aes()}@*)
  geom_histogram(position = "dodge") +           (*@\ixch{geom\_histogram()}@*)
  theme_bw() + xlab("Probability of Yes") +       (*@\ixch{theme\_bw()}\ixch{xlab()}@*)
  scale_fill_manual(values = c("blue", "orange")) (*@\ixch{scale\_fill\_manual()}@*)
# same plot with geom_boxplot() in place of geom_histogram(position = "dodge")  (*@\ixch{geom\_boxplot()}@*)
\end{rcode}
\figref{27}
\fig{Model A horizontal \ix{boxplot}s (y-axis $-0.2$--0.2 meaningless: no y aesthetic). No: box 0.08--0.25, median 0.12,
whiskers 0--0.49, outliers 0.51--0.81. Yes: box 0.54--0.93, median 0.79, whiskers 0.04--1.0. Boxes do not overlap.}

\subsection{Model B: ROC curve}
For poorly predicting models the threshold\ixh{threshold}{} choice barely matters --- they NEVER reach TPR\ixh{TPR}{} 1 with
FPR\ixh{FPR}{} 0 (a ``perfect'' classifier can).
\figref{26}
\begin{routput}
# A tibble: 9 x 3
  .threshold specificity sensitivity      (*@\ixch{.threshold}@*)
1   -Inf          0            1
2      0.113      0            1
3      0.193      0.0182       1
4      0.228      0.0364       1
5      0.267      0.0545       1
6      0.311      0.0788       1
7      0.358      0.0909       0.951
8      0.408      0.285        0.689
9    Inf          1            0
\end{routput}
\fig{Model B \ix{histogram} of \rc{.pred\_Yes} (x 0.1--0.42, y 0--120): few distinct values (one per \rc{sib\_sp}):
0.113, 0.193, 0.228: No 3 each; 0.267: No 4; 0.311: Yes 5, No 2; 0.358: Yes 27, No 32; 0.408: Yes 71, No 118.
Classes at the same values --- no separation; nothing above 0.41.}
Thresholds = grid over the observed predicted probabilities. No predicted P(Yes) $>0.4076525$ = the largest \ix{threshold}
in its ROC curve\ixh{ROC curve}; larger thresholds would give the same results as \rc{Inf}. \ix{boxplot} shows the overlap:
\fig{Model B boxplots: No and Yes essentially identical: box 0.358--0.408, median at upper edge (0.408), left whisker
0.31; No also has outliers at 0.113, 0.193, 0.228, 0.267. Complete overlap.}

\subsection{Conclusions}
\ixh{model evaluation}%
\ix{ROC curve} can also pick the optimal threshold(s)\ixh{threshold}{} for a chosen model --- less common, since defaults
usually balance \ix{sensitivity} and \ix{specificity} well; but to be extra careful to predict \textbf{every} 0 (or every 1)
correctly, shift the threshold that way. Multiclass\ixh{multiclass classification}{} outcome ($>2$ values): one curve per
class for how well the model distinguishes \textbf{each} class (examples in homework/labs).
